\documentclass{exam}

\newif\ifprintselected   % required by exam_config's \ansbox

%%%%%%%%%%%%%%%%%%%%%%%%%%%%%%%%%%%%%%%%%%%%%%%%%%%%%%
%%% SET THESE IF VARIABLES
%%%%%%%%%%%%%%%%%%%%%%%%%%%%%%%%%%%%%%%%%%%%%%%%%%%%%%
% Term, for the running header.
\ifdefined\thisTermIn
	\def\thisterm{\thisTermIn}
\else
	\def\thisterm{Fall 2026}
\fi
\def\thisclass{ECE 444}

% The design frequency. Everyone builds the same antenna here.
\def\fOne{915}

% The second-frequency window. Cadets choose f2 inside [\fTwoLo, \fTwoHi] MHz
% and at least \fTwoGuard\% away from \fOne, which puts the dipole between
% roughly 0.42 and 0.55 wavelengths -- detuned enough that the match is real
% work, not so far that a two-element L-network cannot reach it with the
% stocked parts.
%
% THE UPPER AND LOWER BOUNDS ARE ALSO RANGE-EQUIPMENT BOUNDS. The window must
% sit inside the usable band of the range's source/reference antenna. Confirm
% that model's low-frequency limit before the project goes out and edit these
% three numbers if it does not reach \fTwoLo.
\def\fTwoLo{800}
\def\fTwoHi{1050}
\def\fTwoGuard{5}
\def\fGuardLo{869}   % \fOne\ less \fTwoGuard\%, rounded
\def\fGuardHi{961}   % \fOne\ plus \fTwoGuard\%, rounded

% Due date and weight. Both are quoted from the syllabus -- change there first.
\def\duedate{2 October, 2359}
\def\projweight{30\%}
\def\midweight{40\%}

%%%%%%%%%%%%%%%%%%%%%%%%%%%%%%%%%%%%%%%%%%%%%%%%%%%%%%
%%% END SET THESE IF VARIABLES
%%%%%%%%%%%%%%%%%%%%%%%%%%%%%%%%%%%%%%%%%%%%%%%%%%%%%%
\input{myPackages_gr}
\usepackage[hidelinks, linkcolor=red]{hyperref}
\input{myShortcuts}
% ECE 444 document fonts.
%
% Body text is Barlow (SIL OFL; the faces and the licence are in
% latex/fonts/Barlow/). Barlow is not in TeX Live, so it is loaded as an
% OpenType face through fontspec -- which means the document must be built
% with lualatex, not pdflatex. build_practice.sh and build_lab.sh both do.
%
% Compiling with pdflatex anyway is not an error: the guard below falls back
% to the LaTeX default so the source still compiles under any engine and can
% be syntax-checked. The resulting PDF just is not the one we ship.
%
% Paths are relative to latex/, which is where both build scripts run
% (each does cd "$(dirname "$0")"). Loading by path rather than by installed
% family name means no system font install and no fontconfig, so a fresh
% clone builds a byte-comparable PDF anywhere.
%
% MATH IS DELIBERATELY LEFT IN COMPUTER MODERN. Barlow has no math companion,
% and CM math under sans body text is already the house pairing: the decks set
% words in Source Sans Pro and let MathJax render symbols in its CM-derived
% face. Moving math to a sans as well would mean unicode-math plus a sans math
% font (Fira Math being the plausible one) and re-checking every equation in
% the course for symbol coverage. That is a separate decision, not a side
% effect of this one.

\usepackage{iftex}
\ifLuaTeX
	\usepackage{fontspec}
	\defaultfontfeatures{Ligatures=TeX}
	\setmainfont{Barlow}[
		Path           = fonts/Barlow/,
		Extension      = .ttf,
		UprightFont    = *-Regular,
		ItalicFont     = *-Italic,
		BoldFont       = *-Bold,
		BoldItalicFont = *-BoldItalic,
	]
	\setsansfont{Barlow}[
		Path           = fonts/Barlow/,
		Extension      = .ttf,
		UprightFont    = *-Regular,
		ItalicFont     = *-Italic,
		BoldFont       = *-SemiBold,
		BoldItalicFont = *-BoldItalic,
	]
	% microtype earns its place here specifically because of Barlow. Its
	% glyphs are narrower than Computer Modern's, so justified prose sets
	% tighter and TeX has to squeeze interword space to fit; two paragraphs
	% in the L02/L03 sets went overfull on the first Barlow build. Protrusion
	% and font expansion absorb that -- both are lualatex features, so this
	% belongs inside the guard with everything else.
	\usepackage{microtype}

	% Mono stays Latin Modern: Barlow ships no monospaced face, and verbatim
	% and code need one that actually is fixed-pitch.
\fi

\setlength{\parindent}{0pt}
\setlength{\parskip}{6pt}

\input{exam_config}
\noprintanswers
\printselectedfalse

\renewcommand{\thesection}{\arabic{section}}
\setcounter{secnumdepth}{2}

% Requirement callout: the numbered things a report must contain. Kept visually
% distinct from prose so a cadet can find every one of them on a second pass.
\newcommand{\req}[1]{%
	\par\vspace{2pt}%
	\noindent\fbox{\begin{minipage}{0.97\linewidth}\vspace{1pt}#1\vspace{1pt}\end{minipage}}%
	\par\vspace{4pt}}

\firstpageheader{}{}{}
\runningheader{\thisclass \hspace{1pt} -- \thisterm }{Midterm Project}{\thepage\hspace{1pt} of\hspace{1pt} \numpages}
\firstpagefooter{}{}{}
\runningfooter{}{}{}
\runningheadrule

\begin{document}
\pagestyle{headandfoot}

%%%%%%%%%%%%%%%%%%%%%%%%%%%%%%%%%%%%%%%%%%%%%%%%%%%%%%
% title block
%%%%%%%%%%%%%%%%%%%%%%%%%%%%%%%%%%%%%%%%%%%%%%%%%%%%%%
\begin{center}
USAF ACADEMY \\
DEPARTMENT OF ELECTRICAL AND COMPUTER ENGINEERING \\
MIDTERM PROJECT (assigned Lesson 9) - Antenna Pattern Measurement \\
\textbf{Project Description} \\
\thisterm \\[4mm]
\end{center}

\fullwidth{\textbf{Due \duedate.} Individual. Worth \projweight\ of the
course grade and \midweight\ of your midterm grade, scored on the
\textbf{100-point rubric in Section 12}. There is no resubmission --- the
grade you earn is the grade.}

\vspace{2pt}

\section{The assignment in one paragraph}

Build a \fOne\ MHz wire dipole. Trim it until it resonates, and prove it did.
Then choose a second frequency, leave the antenna exactly as it is, and build
a matching network that presents 50 \ohms at that new frequency
instead. Model the antenna in NEC at both frequencies. Take it into the
chamber and measure its pattern at both frequencies, co-polarized and
cross-polarized, against a reference antenna of known gain. Then write up what
you predicted, what you measured, and what you make of the difference.

You are not designing an antenna. Lesson 7 gave you the design and Lesson 8
let you model it, so the pattern, the gain, and the impedance are all
predictable before you measure anything. The work of this project is the
comparison.

\section{The engineering question}

There are two ways to make an antenna present 50 \ohms at the
connector, and this project makes you do both to the same piece of wire.

At \fOne\ MHz you change the \emph{antenna}: you cut it until its own
reactance vanishes. At your second frequency $f_2$ you change the
\emph{network}: the antenna stays wrong and a pair of components hides it from
the source. Both roads end at VSWR $\le$ 2 looking into the connector. They do
not end at the same place anywhere else.

\req{\textbf{Your report answers this question:} what did each method cost
you, and what did the radiated pattern do about it? A VSWR reading at the
connector is a statement about power accepted, not about power radiated or
where it went. Your pattern measurements are how you tell the difference.}

That question is the spine of the report. Everything below exists to give you
the evidence to answer it.

\section{What this project assesses}

The project is where the Module 2 measurement objectives are demonstrated end
to end:

\begin{itemize}[nosep]
	\item I can set up a one-port sweep, run an SOL calibration at the correct
	      reference plane, and verify the calibration before I trust a reading.
	\item I can read resonance, impedance, and VSWR $\le$ 2 bandwidth off a
	      trace and off the Smith chart, and say from the sign of the reactance
	      which way to trim.
	\item I can verify a range geometry is valid for both antennas before I
	      take data, using all three far-field criteria.
	\item I can acquire principal-plane cuts and normalize, plot, and annotate
	      them in dB down from the peak.
	\item I can extract half-power beamwidth, first sidelobe level,
	      front-to-back ratio, gain by comparison, and cross-polarization
	      discrimination from measured data.
	\item I can measure my own noise floor, state the resulting dynamic range,
	      and say which of my numbers clears it and by how much.
\end{itemize}

\section{Phase 1 - Build the dipole}

Predict before you cut. Lesson 7 gives the half-wave dipole's dimensions, its
$73 + j42.5$ \ohms\ input impedance, and its 2.15 dBi directivity. Work out
the numbers for \fOne\ MHz on paper first.

Build it from wire and an SMA connector. The example antenna on the bench
shows the feed construction: arms soldered to the center pin and to the body,
a clean symmetric junction, arms straight and colinear. Make yours mechanically
rigid --- an arm that moves when you move the coax makes every later
measurement a different antenna.

\req{\textbf{Required:} the predicted arm length and predicted $Z_{in}$ at
\fOne\ MHz, computed \emph{before} the first cut, quoted in the report next to
what you actually built.}

Cut long. You can remove wire; you cannot put it back.

\section{Phase 2 - Find and fix the resonance}

Use the handheld VNA for this phase. It lives at the bench, it sweeps fast,
and trimming is an iterative loop you will run several times.

\begin{enumerate}[nosep]
	\item Run a short-open-load calibration and verify it. State your reference
	      plane --- the numbers mean nothing without it, and the length of cable
	      past that plane is pure rotation on the chart.
	\item Sweep the antenna. Find the frequency where the reactance crosses
	      zero. That is the resonance, and it is almost certainly not \fOne\
	      MHz on the first try.
	\item Read the sign of the reactance at \fOne\ MHz and let it tell you
	      which way to trim.
	\item Trim, re-measure, repeat.
\end{enumerate}

\textbf{Log every iteration.} Arm length, resonant frequency, impedance at
resonance, VSWR at \fOne\ MHz. You need four or five of these rows to make the
trim history plot, and the plot is worth more than the final number: it shows
a physical trend you can compare against the $\lambda/2$ prediction.

\req{\textbf{The bar:} VSWR $\le$ 2 at \fOne\ MHz, measured, with the trace
shown. Switch the analyzer to VSWR format at the bench so the limit line is
simply 2.}

Note what the textbook antenna does here. An ideal $\lambda/2$ dipole at
$73 + j42.5$ \ohms\ on a 50 \ohms\ system is VSWR 2.18 --- it
\emph{fails} this bar. Trimming to resonance is what passes it. Say in your
report what your trimmed antenna's resistance came out to be and whether that
is consistent with the 73 \ohms\ Lesson 7 derived.

\section{Phase 3 - Build a balun}

A coaxial cable is unbalanced and a dipole is balanced. Connect them directly
and current that should have stayed on the inside of the shield flows back
down its outside, which radiates. The thing you then measure in the chamber is
not your dipole. It is your dipole plus a length of coax, and it will show up
as a pattern that is tilted, asymmetric, and dependent on how the cable is
routed that day.

You will build a balun. Two constructions work at \fOne\ MHz; pick one.

\textbf{Quarter-wave sleeve balun.} A conducting sleeve (copper tape over a
tube, or thin brass tubing) coaxial with the feed line, shorted to the shield
at the end away from the antenna and open at the feed end. Length is a quarter
wavelength in air at \fOne\ MHz --- compute it. The shorted quarter wave
presents a high impedance at the feed, which is what chokes the outside
current.

\textbf{Coiled-coax choke.} Wind several turns of the feed coax into a tight
solenoid a few centimeters across, close to the feed point. The coil's
inductance and its self-capacitance form a high impedance to common-mode
current. It is quicker to build than a sleeve and its performance is less
predictable, which is why the verification below matters.

\req{\textbf{Required evidence that it works.} Measure the antenna's VSWR at
\fOne\ MHz while you move, coil, and grip the feed line, with and without the
balun. Report how much the reading swings in each case. A working balun makes
the antenna's behavior stop depending on its cable. If you have chamber time
for it, the stronger evidence is a pattern cut taken both ways.}

Do not skip this and do not do it last. The balun is on the antenna for every
measurement that follows, including the trim verification you may need to
repeat once it is fitted.

\section{Phase 4 - Detune, predict, and match}

Now the second half of the engineering question.

\subsection{Choose $f_2$}

Pick any $f_2$ between \fTwoLo\ and \fTwoHi\ MHz that is at least
\fTwoGuard\% away from \fOne\ MHz --- that is, outside \fGuardLo--\fGuardHi\
MHz. Different cadets should pick different frequencies. Declare your choice
at the Lesson 14 checkpoint.

\subsection{Predict $Z_a(f_2)$ before you design anything}

Model your antenna --- the one you actually built, at the length you actually
trimmed it to --- in the course NEC simulator and read its input impedance at
$f_2$. This prediction is the load your matching network is designed against.

This ordering is not a formality. A matching network designed against a
measured impedance is a fitting exercise; one designed against a prediction
and then measured is an experiment. Do it in that order and report both
numbers.

\subsection{Design, build, verify}

Use the L-network method from the Lesson 4 lab. Your network must absorb the
antenna's own reactance at $f_2$ as well as transform its resistance, so work
from the full complex $Z_a(f_2)$, not from its real part. Build it on the SMA
proto board with discrete components; choose the nearest stocked values and
report the difference between the values you designed and the values you
soldered.

Two elements should reach any $f_2$ in the allowed window. If yours will not,
say why in the report before adding a third.

Verify on the Copper Mountain VNA. This is your final, reportable impedance
data, so calibrate carefully and state the reference plane again.

\req{\textbf{The bar:} VSWR $\le$ 2 at $f_2$, measured on the Copper Mountain,
with the trace shown and the VSWR $\le$ 2 bandwidth quoted. Compare that
bandwidth against the resonant antenna's bandwidth at \fOne\ MHz. They will
not be the same, and the difference is one of the costs the engineering
question is asking about.}

\subsection{Measure the same antenna on both instruments}

Once the antenna is final, sweep it on \emph{both} the handheld and the Copper
Mountain, calibrated to the same reference plane, over the same span. Overlay
the two traces.

They will not lie on top of each other. Report the disagreement --- where it is
worst, how large it is, and what you think accounts for it. Two instruments
disagreeing about one unchanged object is the most honest uncertainty estimate
you will get in this course, and it is worth more than any figure you could
quote from a data sheet.

\section{Phase 5 - Simulate}

Use the course NEC2 simulator. Model the antenna you built: your trimmed
length, your wire diameter, your feed geometry.

Produce, at both \fOne\ MHz and $f_2$:

\begin{itemize}[nosep]
	\item input impedance and VSWR referenced to 50 $\Omega$,
	\item E-plane and H-plane pattern cuts,
	\item directivity, and the gain if you model loss.
\end{itemize}

State your model's assumptions where they differ from the hardware: perfect
conductor, no balun, no connector, no feed cable, free space rather than a
chamber. Every one of those is a candidate explanation for a later
disagreement, and a disagreement you predicted is worth far more than one you
discovered.

\section{Phase 6 - Measure the pattern}

\subsection{Before you take data}

Verify the range geometry using all three far-field criteria from Lesson 9,
for \emph{both} the antenna under test and the reference antenna, at
\emph{both} frequencies. State which criterion governs. A dipole is small
compared to a wavelength, so $2D^2/\lambda$ is probably not the one that binds
--- find out which is.

\subsection{Measure your floor}

Before the pattern, measure the floor: with the source off, or with the
antennas cross-polarized and rotated away, record the level the receiver
reports. Quote it in dB below your co-polarized peak. That number is your
dynamic range, and it bounds every quantity you extract afterward.

\req{\textbf{Draw the floor as a horizontal line on every pattern plot in your
report.} A pattern measurement has a floor; the chamber contributes one and
the receiver contributes the other, and whichever is higher is the one you
have. Any number that lands within a few dB of it is a \emph{bound}, not a
value, and must be quoted as ``$\le$'' in your tables.}

\subsection{The cuts}

At both \fOne\ MHz and $f_2$:

\begin{itemize}[nosep]
	\item E-plane cut, co-polarized,
	\item H-plane cut, co-polarized,
	\item both planes again, cross-polarized.
\end{itemize}

Keep the range, the cabling, and the source power identical between the co-
and cross-polarized runs. If you change anything, you have changed the
reference the two are compared against and the discrimination number is
meaningless.

\subsection{The reference antenna}

Measure the reference antenna's own pattern as well. It is a five-minute swap
and a five-minute sweep, and it buys you two things: a check that your range
is behaving, since you know what that antenna's pattern should look like, and
the received-power reference for gain by comparison.

Gain by substitution: replace your antenna with the reference, change nothing
else, and record the received power. Then
%
\eq{
	G_{\text{AUT}} \lb \dbi \rb = G_{\text{ref}} \lb \dbi \rb
	+ \lp P_{\text{AUT}} - P_{\text{ref}} \rp \lb \db \rb.
}
%
Note the conditions this is valid under --- Lesson 9 lists them --- and say
whether yours were met.

\section{Phase 7 - Analysis}

This is the part of the report that carries the grade.

\subsection{Normalize honestly}

The chamber gives you relative received power. NEC gives you absolute
directivity. They are not the same quantity and you cannot simply plot them on
one axis.

\req{\textbf{State your normalization basis in the caption of every overlay.}
The default: normalize the measured co-polarized cut to its own peak,
normalize the simulated cut to its own peak, and plot both in dB down from
peak. That comparison is about pattern \emph{shape} and it is honest. Absolute
level is a separate claim, made separately, from the gain substitution
measurement.}

What you must not do is slide one curve against the other until they agree.
Shape agreement and gain agreement are two different findings, and collapsing
them into one plot hides whichever of them failed.

\subsection{Extract the numbers}

From the measured cuts, at both frequencies: half-power beamwidth in both
planes, null depths and positions, first sidelobe level in dBc,
front-to-back ratio, cross-polarization discrimination, and gain. For each,
say whether it clears your floor.

\subsection{Reconcile}

Put prediction and measurement side by side and account for the differences.
Lesson 7's closed form, your NEC model, and your chamber data are three
independent statements about one antenna. Where they disagree, work out which
one you believe and why.

Candidate explanations you should be able to rule in or out: the feed cable
and balun, chamber reflections and the quiet-zone limits, the range geometry,
alignment and polarization error, the model's idealizations, connector and
network loss at $f_2$, and your own repeatability.

\req{\textbf{This is where the grade is won.} You are not required to make the
numbers agree. You are required to know why they do not. A disagreement you
have diagnosed is a result; a disagreement you have not mentioned is a
missing result.}

There is no error budget to fill in. The errors are there whether or not you
tabulate them --- what the report needs is your reasoning about which of them
is responsible for what you actually see.

\section{The report}

One typed technical report, submitted as a PDF, plus your data files. There is
no page limit and no template: the structure is yours to choose, and choosing
it is part of the work. Two pages will not do it.

\textbf{Nothing handwritten.} Every figure is generated by software. Every
axis is labeled with units. Every trace is in a legend. Every figure and table
is numbered and has a caption that says what the reader should take from it.
An uncaptioned instrument screenshot is not a figure.

\subsection{Required figures}

\begin{enumerate}[nosep]
	\item Annotated photograph of the antenna as built, with dimensions and the
	      balun visible.
	\item Trim history: VSWR against frequency for each cut, on one set of
	      axes, with the VSWR $=$ 2 line drawn.
	\item Smith chart at $f_2$: the antenna unmatched and the antenna matched.
	\item Matching network schematic, with designed and as-built component
	      values.
	\item Handheld and Copper Mountain traces of the final antenna, overlaid,
	      same reference plane.
	\item The NEC model geometry.
	\item Simulated E- and H-plane patterns at both frequencies.
	\item Measured co- and cross-polarized cuts, both planes, both
	      frequencies, with the measured floor drawn on each.
	\item Measured against simulated, overlaid, normalization basis stated in
	      the caption.
	\item The reference antenna's measured pattern.
\end{enumerate}

\subsection{Required tables}

\begin{enumerate}[nosep]
	\item \textbf{Predicted against measured.} Resonant frequency, input
	      impedance, VSWR, VSWR $\le$ 2 bandwidth, half-power beamwidth in both
	      planes, first sidelobe level, front-to-back ratio, cross-polarization
	      discrimination, and gain --- with three columns: Lesson 7 closed
	      form, NEC, and measured. Both frequencies.
	\item \textbf{The match at $f_2$.} Predicted $Z_a(f_2)$, designed component
	      values, as-built component values, achieved VSWR, achieved bandwidth.
\end{enumerate}

\subsection{Data files}

Submit alongside the PDF, named so a reader can tell what each one is:
Touchstone \texttt{.s1p} files from both analyzers, your pattern data as CSV,
and your NEC input decks. Your data files must be your own measurements.

\section{Grading}

The project is scored out of 100 on the rubric below. Each row is worth the
points shown, awarded on how completely and how well the work is done --- not
as all-or-nothing. There is no resubmission, so read this section before you
start, not after.

\begin{center}
\begin{tabular}{|p{4.6in}|c|}
\hline
\textbf{What is assessed} & \textbf{Points} \\
\hline
\textbf{Analysis and reconciliation.} Every disagreement between predicted and
measured either explained or explicitly owned as unexplained, with the
candidate causes ruled in or out. & 20 \\
\hline
\textbf{The two matches.} VSWR $\le$ 2 at \fOne\ MHz by trimming and at $f_2$
by matching, both with traces shown --- or, if a bar was missed, a quantitative
diagnosis of why. & 15 \\
\hline
\textbf{Pattern measurements.} Co- and cross-polarized cuts in both planes at
both frequencies, with the measured floor stated as a number and drawn on the
plots. & 15 \\
\hline
\textbf{Prediction before design.} The NEC prediction of $Z_a(f_2)$ presented,
and presented as what the match was designed against. & 10 \\
\hline
\textbf{Gain by comparison.} The substitution measurement, the reference
antenna's own pattern, and the validity conditions addressed. & 10 \\
\hline
\textbf{Report quality.} All required figures and tables present, typed,
legible, captioned, axes labeled, with the data files. & 10 \\
\hline
\textbf{The balun.} Built, and evidence in the report that it does
something. & 8 \\
\hline
\textbf{Instrument comparison.} Both analyzers' traces of the final antenna,
overlaid at the same reference plane, with the disagreement discussed. & 7 \\
\hline
\textbf{Honest normalization.} Every overlay's comparison basis declared in
its caption. & 5 \\
\hline
\textbf{Total} & \textbf{100} \\
\hline
\end{tabular}
\end{center}

Two things are worth saying plainly about that table. \textbf{Analysis is the
largest single row}, larger than either measurement row, because a campaign
you cannot interpret is not a result. And a report in which everything agrees,
nothing is questioned, and no measurement appears to have surprised its author
will lose most of those twenty points however clean the data is.

Work that is missing entirely earns nothing for that row. Work that is present
but thin --- a pattern plotted without its floor, a match verified but not
explained --- earns part of it.

\section{Checkpoints}

These are not graded. They exist because the chamber is a shared resource and
because a problem found at Lesson 11 is an afternoon, while the same problem
found on 1 October is a failure.

\begin{center}
\begin{tabular}{|l|p{3.9in}|}
\hline
\textbf{By} & \textbf{You should have} \\
\hline
Lesson 11 & A trimmed dipole with a balun on it, VSWR $\le$ 2 at \fOne\ MHz,
shown to the instructor. \\
\hline
Lesson 14 & $f_2$ declared, $Z_a(f_2)$ predicted in NEC, the L-network
designed on paper. \\
\hline
Lesson 15 & Matching network built and verified on the Copper Mountain.
Chamber time booked. \\
\hline
Lesson 17 & Chamber measurements complete, including the reference antenna. \\
\hline
\duedate & Report and data files submitted. \\
\hline
\end{tabular}
\end{center}

Chamber time is scheduled in pairs. \textbf{You may share a chamber session
and help each other take data. The data reduction, the analysis, and the
report are individual.} Two reports containing the same sentences are two
reports that are not your own work.

\section{Authorized resources}

Authorized without restriction: the course site and every lesson, lab, and
practice set in it; your own notes; the course textbook; the NEC2 simulator;
Python and its scientific libraries for reduction and plotting; the
instrument manuals; and your instructor.

Discussion with classmates about method is encouraged. Sharing measured data,
figures, or written text is not.

\textbf{Generative AI use on this project is Level 3.} If you use it, include
a documentation statement, as the syllabus requires on any assignment --- name
the tool, what you used it for, and where it appears in your work.

The litmus test from the syllabus applies here as it does everywhere else: you
must be able to explain every step of what you submit.

\vspace{6pt}
\hrule
\vspace{6pt}

\begin{center}
\textit{Questions about scope, equipment, or scheduling go to the instructor
early. A question asked at Lesson 12 is free.}
\end{center}

\end{document}
